\bookchapter{Глава 50: Быстрый старт с Kotlin — первый язык Android-разработки}

\chapterintro{В этой главе мы наконец учим язык. Не волнуйтесь — Kotlin славится дружелюбием к новичкам, и учить грамматику наизусть не нужно: освоив 20\% синтаксиса, вы будете понимать и писать 80\% повседневного кода. Остальное подсмотрите по ходу проектов. Помните наш принцип обучения: понимать важнее, чем запоминать; уметь пользоваться важнее, чем владеть в совершенстве.}

\booksection{Переменные и типы: даём данным имена}

Программа — это «данные + логика». Переменная — это имя, под которым данные хранятся. В Kotlin два способа объявления:

\begin{lstlisting}[language=Kotlin]
val name = Иван        // val: переменная только для чтения (по умолчанию берём её)
var age = 18             // var: изменяемая переменная (можно менять позже)

age = 19                 // <i class="fas fa-check-circle mr-1"></i> var можно изменить
// name = "Мария"         // <i class="fas fa-times-circle mr-1"></i> val изменить нельзя — ошибка компиляции

val score: Int = 100     // тип можно указать и явно
val pi = 3.14            // можно не писать — выведется сам
\end{lstlisting}

Часто используемых типов всего несколько — запомните их одной таблицей:

\begin{center}
\begin{tabularx}{\linewidth}{|X|X|X|}
\hline
\textbf{Тип} & \textbf{Что хранит} & \textbf{Пример} \tabularnewline
\hline
Int & Целое число & 42 \tabularnewline
Double & Дробное число & 3.14 \tabularnewline
String & Текст & "Привет" \tabularnewline
Boolean & Истина / Ложь & true / false \tabularnewline
List< T> & Набор данных & listOf(1, 2, 3) \tabularnewline
\hline
\end{tabularx}
\end{center}

\begin{center}\small\color{gray}Базовые типы Kotlin, используемые чаще всего\end{center}

\begin{tcbinfo}{Совет: принцип «val прежде всего»}
Железное правило сообщества Kotlin: по умолчанию val, и только при действительной необходимости var. Если там, где хватит val, вы используете var, IDE подсветит жёлтым. Эта привычка избавит вас от целого класса багов «переменную изменили случайно» — чем меньше в коде изменчивости, тем легче о нём рассуждать.

\end{tcbinfo}

\booksection{Функции: упаковываем логику в «кнопки»}

Функция — это именованный фрагмент логики, как кнопка на пульте: нажали (вызвали) — выполнилась цепочка действий. В Kotlin функция объявляется словом fun:

\begin{lstlisting}[language=Kotlin]
fun greet(name: String): String {
    return "Привет, $name!"     // знак $ вставляет переменную в строку
}

println(greet(Иван))        // вывод: Привет, Иван!

// однострочную функцию можно записать выражением
fun add(a: Int, b: Int) = a + b

println(add(3, 5))            // вывод: 8
\end{lstlisting}

Несколько моментов: параметры пишутся в скобках (сначала имя, затем тип, через двоеточие); тип возвращаемого значения — после скобок; строковые шаблоны \$имя\_переменной — фирменное удобство Kotlin: больше не нужно, как в Java, склеивать «строка + переменная».

\booksection{Условия и циклы: учим программу «думать»}

Ум программы складывается из двух вещей: умения делать выбор (if/when) и умения повторять (for/while).

\begin{lstlisting}[language=Kotlin]
// Ветвление: if / else
val age = 17
if (age >= 18) {
    println("взрослый")
} else {
    println("несовершеннолетний")
}

// Множественный выбор: when (в сто раз удобнее switch)
val score = 85
val level = when {
    score >= 90 -> "отлично"
    score >= 60 -> "зачёт"
    else -> "Так держать!"
}

// Повторение: цикл for
val fruits = listOf("яблоко", "банан", "апельсин")
for (fruit in fruits) {
    println("Я люблю $fruit")
}
\end{lstlisting}

when — самый ходовой в Kotlin способ записать множественный выбор; в Android-разработке вы будете пользоваться им каждый день (например, чтобы показывать разные экраны в зависимости от состояния).

\booksection{Null-безопасность: «спасательный круг» Kotlin}

Самый знаменитый баг эпохи Java — NullPointerException (исключение нулевого указателя): программа обращается к тому, в чём «ничего нет», и мгновенно падает. По статистике он годами держит первое место в рейтингах сбоев Android. Kotlin решил проблему на уровне языка: null по умолчанию запрещён.

\begin{lstlisting}[language=Kotlin]
var name: String = "Иван"   // обычный тип: null не допускается
// name = null                // <i class="fas fa-times-circle mr-1"></i> ошибка компиляции — вас сразу остановят

var nickname: String? = null // добавляем вопросик: null разрешён
nickname = "Ваня"

// С nullable-типами нужно обращаться безопасно:
println(nickname?.length)    // безопасный вызов: вернёт null, не падая
val len = nickname?.length ?: 0   // оператор Элвиса: значение по умолчанию при null
println(len)
\end{lstlisting}

\begin{tcbwarn}{Внимание: не путайте ? и ?}
String? означает «может быть пустым», а a ?: b — «если a пусто, возьми b» (запись похожа на причёску Элвиса, поэтому оператор прозвали Elvis). Когда в начале пути компилятор ругается на «ошибки, связанные с null», в восьми случаях из десяти вы забыли обработать nullable-тип — язык просто защищает вас.

\end{tcbwarn}

\booksection{Классы и объекты: формы для изготовления «деталей»}

Класс — это «форма» для объектов, объект — «готовая деталь», отлитая по форме. Всё в Android — экраны, кнопки, данные — объекты. В Kotlin есть волшебный инструмент data class (класс данных), созданный специально, чтобы «носить данные»: одна строка заменяет десятки строк Java:

\begin{lstlisting}[language=Kotlin]
// Обычный класс: форма
class Dog(val name: String) {
    fun bark() = println("$name: Гав!")
}

val dog = Dog("Шарик")      // создаём объект
dog.bark()                 // Шарик: Гав!

// data class: коробка для данных (с готовыми equals/copy и т.д.)
data class User(val id: Int, val name: String, val email: String)

val me = User(1, Иван, "zhangsan@example.com")
val twin = me.copy(name = "Иванушка")   // копируем и меняем одно поле
\end{lstlisting}

Когда позже вы будете делать учёт расходов, каждая запись станет data class: data class Record(val amount: Double, val category: String, val date: String). Смотрите — вы уже способны писать структуру кода настоящего проекта.

\booksection{Лямбды и функции высшего порядка: «синтаксический сахар» Kotlin}

В Android-коде вы ежедневно будете встречать «функции без имени» вида \{ ... \} — они называются лямбдами (анонимными функциями). Button с onClick = \{ count++ \} — это лямбда: мы упаковываем «что делать после нажатия» в один блок и вручаем его кнопке.

\begin{lstlisting}[language=Kotlin]
// Обычный способ: определяем именованную функцию и передаём
fun double(x: Int) = x * 2
val result = listOf(1, 2, 3).map(::double)

// Через lambda: определяем на месте, коротко и прямо
val result2 = listOf(1, 2, 3).map { it * 2 }   // it = текущий элемент

// Комбо «фильтр + преобразование» (крайне часто в реальных проектах)
val bills = listOf(12.5, 88.0, 5.0, 200.0)
val bigBills = bills.filter { it > 50 }.sorted()
// [88.0, 200.0] — «счета больше 50, отсортированы»
\end{lstlisting}

Запомните несколько частых «комбо» для коллекций: filter (отбор), map (преобразование), sorted (сортировка), sum (сумма). 90\% будущей работы со списками — комбинации этих четырёх приёмов. Философия Kotlin: если логику можно высказать одной фразой, не пишите пять.

\booksection{Шпаргалка по синтаксису для тех, кто с нуля}

Не пытайтесь выучить весь синтаксис за один день. Сохраните эту карточку и заглядывайте в неё, пока пишете код, — недели через две всё станет привычным само собой:

\begin{center}
\begin{tabularx}{\linewidth}{|X|X|}
\hline
\textbf{Что нужно сделать} & \textbf{Как записать} \tabularnewline
\hline
Объявить переменную только для чтения & \texttt{val x = 1} \tabularnewline
Объявить изменяемую переменную & \texttt{var x = 1} \tabularnewline
Объявить функцию & fun имя(параметры): тип Результата \{ ... \} \tabularnewline
Условие & if (условие) \{...\} else \{...\} \tabularnewline
Множественный выбор & when \{ условие -> результат \} \tabularnewline
Перебор & for (x in список) \{ ... \} \tabularnewline
Склейка строк & "Привет, \$name" \tabularnewline
Обработка nullable & \texttt{x?.\allowbreak length ?: 0} \tabularnewline
Объявить класс данных & \texttt{data class X(val a: Int, val b: String)} \tabularnewline
\hline
\end{tabularx}
\end{center}

\begin{center}\small\color{gray}Шпаргалка по синтаксису Kotlin (версия для новичков)\end{center}

\begin{tcbpractice}{Практика: поиграйте 10 минут в Kotlin, ничего не устанавливая}
Откройте в браузере play.kotlinlang.org (официальную онлайн-песочницу Kotlin), вставьте туда каждый пример кода из этой главы, запустите, поменяйте числа и текст и посмотрите, как меняется вывод. Эти 10 минут «сначала поиграть, потом учить» вдвое облегчат вам старт в главе 51.

\end{tcbpractice}

\begin{tcbsummary}{Коротко о главном}
\begin{itemize}
  \item Предпочитайте val, к var прибегайте с осторожностью; тип можно вывести автоматически, но можно объявить и явно после двоеточия
  \item Функция = упакованная логика; строковый шаблон \$ — фирменное удобство
  \item if/when — выбор, for — цикл; when в Android-разработке «мастер на все руки»
  \item Null-безопасность — спасательный круг Kotlin: обычные типы null отвергают, nullable-типы требуют обработки через ? и ?:
  \item data class — коробка для данных: одно поле на строку, в настоящих проектах используется ежедневно
  \item 20\% синтаксиса делают 80\% работы: сначала научитесь пользоваться, а мастерство придёт потом
\end{itemize}

\end{tcbsummary}
